\documentclass[12pt]{exam}
\usepackage{amsthm}
\usepackage{bm}
\usepackage{libertine}
\usepackage[utf8]{inputenc}
\usepackage[margin=0.5in]{geometry}
\usepackage{amsmath,amssymb}
\usepackage{multicol}
\usepackage[shortlabels]{enumitem}
\usepackage{siunitx}
\usepackage{physics}
\usepackage{booktabs}
\usepackage{graphicx}
\usepackage{pgfplots}
\usepackage{listings}
\usepackage{tikz}
\usepackage{tikz-3dplot}

\bmdefine{\ii}{i}
\bmdefine{\jj}{j}
\bmdefine{\kk}{k}
\newcommand{\te}{\theta}
\newcommand{\cC}{\mathcal{C}}
\newcommand{\word}[1]{\quad\text{#1}\quad}

\pgfplotsset{width=10cm,compat=1.9}
\usepgfplotslibrary{external}
\tikzexternalize

\newcommand{\class}{Math 261-001}
\newcommand{\examnum}{Quiz 3 Solutions}
\newcommand{\examdate}{September 11}

\begin{document}
\pagestyle{plain}
\thispagestyle{empty}

\noindent
\textbf{\class}\\
\textbf{\examnum}, \textbf{\examdate}

\vspace{10pt}

\small
\begin{enumerate}

\item \textbf{Final answer:}
\[
  \boxed{-1}.
\]
Holding the indicated variable constant, the ideal-gas equation gives
\[
  \left(\frac{\partial V}{\partial T}\right)=\frac{nR}{P},
  \qquad
  \left(\frac{\partial T}{\partial P}\right)=\frac{V}{nR},
  \qquad
  \left(\frac{\partial P}{\partial V}\right)
  =-\frac{nRT}{V^2}.
\]
Therefore,
\[
  \frac{nR}{P}\cdot\frac{V}{nR}\cdot
  \left(-\frac{nRT}{V^2}\right)
  =-\frac{nRT}{PV}=-1,
\]
because $PV=nRT$.

\item \textbf{Final answer:}
\[
  \boxed{
  \frac{\partial^2z}{\partial x^2}
  =6yf_u+36x^2y^2f_{uu}+48xy^3f_{uv}+16y^4f_{vv}}.
\]
All derivatives of $f$ in this expression are evaluated at
\[
  (u,v)=(3x^2y,4xy^2).
\]
First,
\[
  u_x=6xy,\qquad v_x=4y^2,
\]
so the chain rule gives
\[
  z_x=f_u u_x+f_v v_x
  =6xyf_u+4y^2f_v.
\]
Differentiating once more with respect to $x$,
\[
  z_{xx}
  =f_{uu}u_x^2+2f_{uv}u_xv_x+f_{vv}v_x^2
   +f_u u_{xx}+f_v v_{xx}.
\]
Here
\[
  u_{xx}=6y,\qquad v_{xx}=0.
\]
Substitution yields
\[
  z_{xx}
  =36x^2y^2f_{uu}+48xy^3f_{uv}+16y^4f_{vv}+6yf_u,
\]
which is the displayed answer.

\newpage

\item \textbf{Final answer:}
\[
  \boxed{\left.\frac{dz}{dt}\right|_{t=1}=2}.
\]
The multivariable chain rule gives
\[
  \frac{dz}{dt}=f_x(x,y)\frac{dx}{dt}
  +f_y(x,y)\frac{dy}{dt}.
\]
At $t=1$, $(x,y)=(2,1)$, $dx/dt=2t=2$, and $dy/dt=2$. Since
$\nabla f(2,1)=(3,-2)$,
\[
  \left.\frac{dz}{dt}\right|_{t=1}=3(2)+(-2)(2)=2.
\]

\item \textbf{Final answer:}
\[
  \boxed{\frac{1}{\sqrt{30}}(-5,-1,-2)}.
\]
The gradient of the temperature is
\[
  \nabla T(x,y,z)
  =\left(18x+6yz,\,-6y+6xz,\,6xy\right).
\]
At the fly's location,
\[
  \nabla T(2,2,2)=(60,12,24)=12(5,1,2).
\]
The direction of fastest decrease is the negative unit gradient. Since
$\lVert(5,1,2)\rVert=\sqrt{30}$, that direction is
\[
  -\frac{\nabla T(2,2,2)}{\lVert\nabla T(2,2,2)\rVert}
  =\frac{1}{\sqrt{30}}(-5,-1,-2).
\]

\end{enumerate}

\end{document}
